\section{Introduction}

Unmanned maritime operations are becoming systems of systems. A UAV can rapidly survey a wide area and maintain a radio link; a USV can loiter, carry compute and communications equipment, and bridge air and underwater segments; a UUV can collect close-range observations where radio and satellite services do not reach. Cross-domain control studies increasingly coordinate such platforms for search, tracking, formation, navigation, and surveillance missions \citep{hu2024coordinated,wei20233u,wu2020cooperative,ke2022cooperative,liu2023formation,ke2023cross,pang2026cooperative}. The complementary capabilities are operationally attractive, but they also couple domains whose communication physics, failure modes, owners, and security assumptions differ substantially. Broad reviews of UAV networks and applications emphasize mobility, limited energy, intermittent links, and safety constraints \citep{hayat2016survey,shakhatreh2019unmanned,mohsan2023unmanned}; reviews of unmanned marine and underwater vehicles add acoustic delay, difficult localization, harsh deployment conditions, and constrained recovery \citep{bae2023survey,wibisono2023survey,gonzlezgarca2020autonomous}. A mission decision may therefore depend on evidence that has crossed several media, gateways, processing stages, and administrative boundaries.

This creates a \emph{mission-trust gap}. Conventional cybersecurity questions---who authenticated, which packet arrived, whether a file hash matches, or whether replicas agree---remain necessary, yet none alone answers whether a particular data product is suitable for a particular decision at a particular time. A correctly signed sonar report can be physically implausible; a valid observation can be replayed after it becomes stale; a correctly ordered ledger entry can originate from a captured endpoint; and a well-provenanced product can still be unsuitable for a high-consequence task because its uncertainty is too large. Research on UAV and swarm security catalogs spoofing, jamming, routing attacks, capture, malware, data manipulation, and privacy threats \citep{mekdad2023survey,wang2025survey,hadi2023comprehensive,yahuza2021internet}. Underwater-network security reviews describe related attacks under low bandwidth, long propagation delay, and expensive retransmission \citep{jiang2019securing,yang2018challenges,yisa2021security,adam2024state}. These literatures identify important mechanisms, but the unit of analysis is often a platform or network rather than the chain from physical observation to mission use.

Blockchain is frequently proposed as an answer to distributed trust. UAV-oriented reviews report uses in identity, authentication, coordination, access control, data sharing, and audit \citep{hafeez2023blockchain,alladi2020applications,mehta2020blockchain,hawashin2024blockchain,phadke2024unifying}. Representative proposals use ledgers for swarm identity, lightweight authentication, trusted message exchange, and distributed data acquisition \citep{han2022identity,tan2022blockchain,liu2023lightweight,islam2019bus,karmakar2024blockchain}. Maritime reviews similarly associate blockchain with documentation, information sharing, authentication, and inter-organizational coordination \citep{farah2024survey,li2021survey,yang2019maritime}. The attraction is understandable: several organizations can maintain a shared, append-oriented record without granting one participant unilateral control over history. However, the word \emph{trust} becomes ambiguous when replicated-state consistency is allowed to stand for sensor truth, endpoint integrity, privacy, or timeliness. General blockchain security and Internet-of-Things (IoT) surveys repeatedly document scalability, resource, governance, privacy, and interoperability constraints \citep{li2020survey,reyna2018blockchain,uddin2021survey,monrat2019survey}. A useful review must therefore say not only where a ledger helps, but what it does not establish and when it is unnecessary.

Adjacent research offers parts of the answer. Zero trust rejects implicit confidence based only on network location and instead emphasizes explicit, least-privilege, continuously evaluated access \citep{rose2020zero,syed2022zero}. Remote attestation separates the production of device evidence, its appraisal, and the relying party's decision \citep{birkholz2023remote,hardjono2021towards,ott2023universal}. Provenance models describe how entities, activities, and agents are related and can support accountability across transformations \citep{w3c2013provdm,w3c2013provo,hu2020survey,pan2023data}. Trust-management research supplies reputation, probabilistic, subjective-logic, fuzzy, and learning-based methods \citep{ahmed2019trust,aaqib2023iot,junior2021survey,jsang2016subjective,mui2003computational,kamvar2003eigentrust}. Yet these mechanisms are normally discussed in separate architectural vocabularies. A cross-domain mission requires them to be aligned around the objects being judged, the assurance claim being made, and the stage at which evidence is needed.

This review advances one central argument: \emph{mission trust is a chain of non-substitutable assurance claims, and blockchain can strengthen selected transitions in that chain but cannot complete the chain by itself}. To make the argument operational, we propose the authors' \OAL{} framework. It relates four trust objects---entity, platform/device, link/path, and data product---to four assurance layers---observation validity, commitment and provenance integrity, ledger and governance consistency, and mission-use suitability---over six lifecycle stages from admission to recovery. A cross-cutting assurance and governance plane carries identities, attestation results, provenance commitments, trust state, policy and model versions, ledger state, and consortium rules. It is a plane of mechanisms, not a fifth mission object.

The review makes five further contributions. First, it represents dependencies through a typed object graph with relations such as \textit{controls}, \textit{hosts}, \textit{observes}, \textit{transmits}, \textit{transforms}, \textit{endorses}, \textit{records}, and \textit{consumes}. This graph joins evidence propagation to attack propagation. Second, it reframes dynamic trust as a context--temporal--graph (CTG) reasoning loop rather than as a new calibrated algorithm: evidence is appraised, attributed, propagated with uncertainty, converted into a policy action, audited, and fed back. Third, it introduces a unified threat-analysis sentence---attacker, capability, object, propagation path, mission impact, and control---while explicitly separating malice, benign environmental degradation, and insufficient evidence. Fourth, it supplies an evidence-maturity ladder from conceptual argument to sustained deployment. Fifth, it proposes a future \CDUTP{} that maps mission-trust evidence to existing zero-trust, attestation, provenance, communications, and permissioned-ledger specifications. The profile is an author proposal, not a current standard.

The scope is deliberately public and defensive. Civil applications and publicly documented military-support settings are included because coalition ownership, contested communication, and mission criticality sharpen the governance problem. The review does not discuss weapon effects, targeting procedures, or sensitive tactics. It reports no new experiment, dataset, performance result, protocol proof, or calibrated trust score. Its originality lies in conceptual integration, boundary setting, and an evaluation agenda that can be tested by future work.

The remainder of the paper describes the review scope and analytical lens (Section~\ref{sec:scope}), establishes the operational and conceptual foundations, compares fragmented research landscapes, develops OAL and CTG reasoning, examines blockchain patterns and boundaries, organizes evaluation evidence by maturity, and closes with research propositions, limitations, and conclusions.
